% Tested many
% 0.5 is best (porous)
% changing f changes scattering properties and a_max
% is not certain


The scattering models for spherical grains at various porosities showed that compact ($f = 1$) and very porous grains ($f = 0.01$) provide significantly poorer fits to the observations compared to the intermediate filling factors. The goodness-of-fit improves as the filling factor decreases from $f = 1$ toward $f = 0.5$, reaches a minimum at $f = 0.5$, and then worsens towards $f = 0.01$. This preference for a porous grain indicates that the scattering properties required to reproduce the observations are best represented by porous grains, rather than a fully a compact or very porous grain. This trend was seen at all fitting cases, suggesting the preferred porosity does not depend on the wavelengths included in the fit. 




% \textbf{Effect on \Pw :}From Figures \ref{fig: Spherical_BEST}, \ref{fig: appendix_all_band}, and \ref{fig: appendix_just_47} we can visually see how changing the porosity ($f$) changes the polarization curves. I will focus on Figure \ref{fig: Spherical_BEST} since that is the best fit. 
% \note{Make \Pw vs a for different porosities, without scale factor. }

% \textbf{Effect on maximum grain size: }From Table \ref{tab: dust_model_fits_spherical}, the maximum best-fit dust grain size increases as the porosity decreases, meaning that when grains are compact, the model best matches the observations at smaller grain sizes, whereas more porous grains need a larger grain size to match the observations. \add{some plot to talk about this more?} \add{model sucess vs grain size}. This shows that the porosity and success of the fit do have some sort of correlation, as the results are not random. 